\subsection{Borda Count Method and Voting Theory Review}

\content{
        \begin{example} Let's just dive right into an example. \\
        \begin{tabular}{|l|c|c|c|c|}
        \hline
        & 51 & 25 & 10 & 14 \\
        \hline
        1st & S&T&P&O\\
        \hline
        2nd & T&P&T&T\\
        \hline
        3rd & O&O&O&P\\
        \hline
        4th & P&S&S&S\\
        \hline
        \end{tabular}
        \end{example}
}{% presentation
\vspace{3in}
}{% nopresentation
\vspace{3in}
}{% comments
Go over how points are assigned, then do the example. $T$ is the winner of this
election, with 325 points.
}

\content{
        \begin{example} Determine who wins under the Borda Count method.\\
        \begin{tabular}{|l|c|c|c|c|c|}
        \hline
        & 44 & 14 & 70 & 22 & 80 \\
        \hline
        1st & G&G&M&M&B\\
        \hline
        2nd & M&B&G&B&M\\
        \hline
        3rd & B&M&B&G&G\\
        \hline
        \end{tabular}
        \end{example}
}{% presentation
\vspace{2in}
}{% nopresentation
\vspace{3in}\newpage
}{% comments
}

\content{
        \begin{example} Your turn! Use the Borda Count method to determine who
        wins the following two elections! \\
        \begin{tabular}{|l|c|c|c|c|}
        \hline
        & 11 & 5 & 10 & 3 \\
        \hline
        1st & A&C&B&A \\
        \hline
        2nd & C&B&C&B\\
        \hline
        3rd & B&A&A&C\\
        \hline
        \end{tabular}
        \end{example}
}{% presentation
\vspace{2in}
}{% nopresentation
\vspace{3in}
}{% comments
}

\content{
        \begin{example} Use the Borda Count method to determine who
        wins the following election! \\
        \begin{tabular}{|l|c|c|c|c|}
        \hline
        &  120 & 50 & 40 & 90 \\
        \hline
        1st & C&B&D&A\\
        \hline
        2nd & D&C&A&C\\
        \hline
        3rd & B&A&B&B\\
        \hline
        4th & A&D&C&D\\
        \hline
        \end{tabular}
        \end{example}
}{% presentation
\vspace{2in}
}{% nopresentation
\vspace{3in}\newpage
}{% comments
}

\content{
        \begin{example} Let's go back to the first example. Recall that the
        winner of the election was $T$ with 325 points. Is there something not
        fair about that? \\
        \begin{tabular}{|l|c|c|c|c|}
        \hline
        & 51 & 25 & 10 & 14 \\
        \hline
        1st & S&T&P&O\\
        \hline
        2nd & T&P&T&T\\
        \hline
        3rd & O&O&O&P\\
        \hline
        4th & P&S&S&S\\
        \hline
        \end{tabular}
        \end{example}
}{% presentation
\vspace{3in}
}{% nopresentation
\vspace{3in}
}{% comments
$S$ has a majority!
}

\content{
        \begin{example} Use the Plurality method to determine the winner of this
        election. \\
        \begin{tabular}{|l|c|c|c|c|}
        \hline
        & 9 & 19 & 11 & 8 \\
        \hline
        1st & B&A&C&B\\
        \hline
        2nd & A&B&B&C\\
        \hline
        3rd & C&C&A&A\\
        \hline
        \end{tabular}
        \end{example}
}{% presentation
\vspace{2in}
}{% nopresentation
\vspace{2in}
}{% comments
}

\content{
        \begin{example} With the same preference schedule as in the previous
        example, who wins under the Pairwise Comparison method? 
        \end{example}
}{% presentation
\vspace{3in}
}{% nopresentation
\vspace{3in}
}{% comments
}

\content{
        \begin{example} Use the Plurality with Elimination method to determine who wins
        this election. \\
        \begin{tabular}{|l|c|c|c|c|c|c|}
        \hline
        & 120 & 50 & 40 & 90 & 60 & 100\\
        \hline
        1st & C&B&D&A&A&D\\
        \hline
        2nd & D&C&A&C&D&B\\
        \hline
        3rd & B&A&B&B&C&A\\
        \hline
        4th & A&D&C&D&B&C\\
        \hline
        \end{tabular}
        \end{example}
}{% presentation
\vspace{3in}
}{% nopresentation
\vspace{3in}
}{% comments
}

\content{
        \begin{example} Time to think! Does the Pairwise Comparison Method
        Satisfy the Majority Criterion? Why or why not? 
        \end{example}
}{% presentation
\vspace{2in}
}{% nopresentation
\vspace{2in}
}{% comments
}

\content{
Fairness Criterion: These are all things that we think should be true if a
voting method is fair. 
\begin{enumerate}
\item Condorcet Criterion: If there is a Condorcet candidate, they should win
the election.
\item Independence of Irrelevant Alternatives (IIA) Criterion: if a loser drops
out of the election, this should not hurt the winners chances of winning.
\item Monotonicity Criterion: If voters change their votes to support the
winner, this should not hurt the winners chances of winning.
\item Majority Criterion: If a candidate has a majority, they should win the
election. 
\end{enumerate}
}{% presentation
}{% nopresentation
}{% comments
}

\content{
\begin{theorem} (Arrow's Impossibility Theorem) No voting method is guaranteed
to always satisfy all four fairness criterion. 
\end{theorem}
}{% presentation
}{% nopresentation
}{% comments
}
